\documentclass[oneside,a4paper,14pt]{extarticle}
\usepackage[a4paper,top=20mm,bottom=20mm,left=30mm,right=10mm]{geometry}
\usepackage{amsmath, amssymb, amsfonts, amsthm}
\usepackage[T2A]{fontenc}
\usepackage[utf8]{inputenc}
\usepackage[russian]{babel}
\usepackage{textcomp}
\usepackage{indentfirst}
\usepackage{graphicx}
\usepackage{float}
\usepackage{wrapfig}
\usepackage{caption}
\usepackage{tocloft}
\renewcommand{\cftsecleader}{\cftdotfill{\cftdotsep}}
\usepackage[breaklinks]{hyperref}
\usepackage{titlesec}
\usepackage{fancyvrb}
\usepackage{xcolor}
\titleformat{\section}{\normalsize\bfseries}{\thesection}{1em}{}
\titleformat{\subsection}{\normalsize\bfseries}{\thesubsection}{1em}{}
\titleformat{\subsubsection}{\normalsize\bfseries}{\thesubsubsection}{1em}{}
\renewcommand\baselinestretch{1.33}\normalsize
\setlength{\parindent}{1.25cm}

\begin{document}

\begin{small}\section*{Сравнительная характеристика программных средств}
	\addcontentsline{toc}{section}{Сравнительная характеристика программных средств}

	В рамках выполнения работы были рассмотрены три основных средства работы с файловой системой: командная оболочка \texttt{cmd}, файловый менеджер Total Commander и графический Проводник Windows. Ниже приведена их подробная сравнительная характеристика.
	\subsection*{Командная строка cmd}

	\textbf{Преимущества:}
	\begin{itemize}
		\item предустановлена изначально в составе операционной системы Windows, не требует дополнительной установки и настройки;
		\item даёт понимание внутренней работы системы, поскольку все операции выполняются через явные команды, что способствует более глубокому изучению архитектуры ОС;
		\item позволяет автоматизировать рутинные операции с помощью пакетных файлов (\texttt{.bat});
		\item минимальное потребление системных ресурсов;
		\item возможность выполнения операций, недоступных или затруднённых в графическом интерфейсе.
	\end{itemize}

	\textbf{Недостатки:}
	\begin{itemize}
		\item неочевидный синтаксис команд по сравнению с оболочками семейства \texttt{bash}, что увеличивает время освоения;
		\item устаревший интерфейс и ограниченные возможности по сравнению с современными командными оболочками (PowerShell);
		\item отсутствие наглядного отображения структуры файловой системы;
		\item отсутствие встроенной подсветки синтаксиса и автодополнения в базовой конфигурации.
	\end{itemize}
	\newpage
	\subsection*{Total Commander}

	\textbf{Преимущества:}
	\begin{itemize}
		\item обладает широким набором полезных функций: двухпанельный режим, встроенный архиватор, FTP-клиент, поиск файлов, массовое переименование;
		\item позволяет существенно повысить скорость работы с файлами за счёт горячих клавиш и групповых операций;
		\item поддержка плагинов и расширений, позволяющая адаптировать программу под конкретные задачи;
		\item удобная работа с большими объёмами файлов и каталогов.
	\end{itemize}

	\textbf{Недостатки:}
	\begin{itemize}
		\item является платным программным продуктом (существует условно-бесплатная версия с напоминанием);
		\item требует ознакомления с обширным функционалом, что увеличивает порог входа для новых пользователей;
		\item не даёт понимания внутренних процессов операционной системы, так как все операции скрыты за графическим интерфейсом.
	\end{itemize}

	\subsection*{Проводник Windows}

	\textbf{Преимущества:}
	\begin{itemize}
		\item предустановлен изначально в составе операционной системы;
		\item понятен большему количеству пользователей за счёт интуитивного графического интерфейса;
		\item интеграция с другими компонентами системы (поиск, библиотеки, облачные сервисы);
		\item поддержка drag-and-drop и контекстного меню.
	\end{itemize}

	\textbf{Недостатки:}
	\begin{itemize}
		\item не даёт понимания внутренних процессов файловой системы;
		\item низкая скорость выполнения массовых операций по сравнению с командной строкой и Total Commander;
		\item ограниченные возможности групповых операций;
		\item устаревший и менее гибкий интерфейс по сравнению с современными файловыми менеджерами.
	\end{itemize}

	% \end{small}\begin{small}\section*{Вывод}
	% \addcontentsline{toc}{section}{Вывод}
	% В ходе выполнения лабораторной работы были изучены и практически освоены основные команды MS DOS для работы с файловой системой: создание, копирование, перемещение, переименование и удаление каталогов и файлов. Были выполнены все шестнадцать пунктов задания, что позволило закрепить навыки работы в режиме командной строки.

	% Кроме того, была проведена сравнительная характеристика трёх программных средств работы с файловой системой: командной строки \texttt{cmd}, файлового менеджера Total Commander и Проводника Windows. В результате анализа установлено, что каждое из средств имеет свои преимущества и недостатки в зависимости от решаемых задач: \texttt{cmd} обеспечивает глубокое понимание работы системы и максимальную гибкость, Total Commander — высокую скорость и расширенный функционал, а Проводник — простоту и доступность для широкого круга пользователей.
\end{small}
\end{document}